\newpage\setcounter{dang}{0}
\section*{ÔN TẬP CUỐI CHƯƠNG}
\subsection{Bài tập trắc nghiệm}
\Opensolutionfile{ans}
\begin{ex}%[0H1Y1-2]
Trong mặt phẳng tọa độ, cặp vec-tơ nào sau đây có cùng phương
\choice
{$\overrightarrow{u}=(2;3)$ và $\overrightarrow{v}=\left(\dfrac{1}{2};6\right)$}
{\True $\overrightarrow{a}=(\sqrt{2};6)$ và $\overrightarrow{b}=(1;3\sqrt{2})$}
{$\overrightarrow{i}=(0;1)$ và $\overrightarrow{j}=(1;0)$}
{$\overrightarrow{c}=(1;3)$ và $\overrightarrow{d}=(2;-6)$}
\loigiai{
Cặp vec-tơ $\overrightarrow{a}=(\sqrt{2};6)$ và $\overrightarrow{b}=(1;3\sqrt{2})$ cùng phương vì $\vec{a}=\sqrt{2}\vec{b}$.
}
\end{ex}
\begin{ex}%[0H1B1-2]
Trong mặt phẳng tọa độ, cặp vectơ nào sau đây vuông góc với nhau?
\choice
{$\overrightarrow{u}=(2;3)$ và $\overrightarrow{v}=(4;6)$}
{$\overrightarrow{a}=(1;-1)$ và $\overrightarrow{b}=(-1;1)$}
{\True $\overrightarrow{z}=(a;b)$ và $\overrightarrow{t}=(-b;a)$}
{$\overrightarrow{n}=(1;1)$ và $\overrightarrow{k}=(2;0)$}
\loigiai{
Cặp vec-tơ $\overrightarrow{z}=(a;b)$ và $\overrightarrow{t}=(-b;a)$ vuông góc với nhau vì $\vec{z}\cdot \vec{t}=a\cdot (-b)+b\cdot a=0$.
}
\end{ex}
\begin{ex}%[0H1Y1-3]
Trong mặt phẳng tọa độ, vectơ nào sau đây có độ dài bằng $1$?
\choice
{$\overrightarrow{a}=(1;1)$}
{$\overrightarrow{b}=(1;-1)$}
{$\overrightarrow{c}=\left(2;\dfrac{1}{2}\right)$}
{\True $\overrightarrow{d}=\left(\dfrac{1}{\sqrt{2}};\dfrac{-1}{\sqrt{2}}\right)$}
\loigiai{
Vec-tơ có độ dài bằng $1$ là vec-tơ $\overrightarrow{d}=\left(\dfrac{1}{\sqrt{2}};\dfrac{-1}{\sqrt{2}}\right)$.\\
Có: $|\vec{d}|=\sqrt{\left(\dfrac{1}{\sqrt{2}}\right)^2+\left(\dfrac{-1}{\sqrt{2}}\right)^2}=1$.
}
\end{ex}
\begin{ex}%[0H1B4-2]
Góc giữa vectơ $\overrightarrow{a}=(1;-1)$ và vectơ $\overrightarrow{b}=(-2;0)$ có số đo bằng
\choice
{$90^\circ$}
{$0^\circ$}
{\True $135^\circ$}
{$45^\circ$}
\loigiai{
Có: $\cos (\vec{a};\vec{b})=\dfrac{\vec{a}\cdot \vec{b}}{|\vec{a}|\cdot |\vec{b}|}=\dfrac{1\cdot (-2)+(-1)\cdot 0}{\sqrt{1^2+(-1)^2}\cdot \sqrt{(-2)^2+0^2}}=\dfrac{-1}{\sqrt{2}}$.\\
Vậy góc giữa hai $\vec{a}$ và $\vec{b}$ bằng $135^\circ$.
}
\end{ex}
\begin{ex}%[0H1B4-2]
Khẳng định nào sau đây là đúng?
\choice
{$(\overrightarrow{a}\cdot\overrightarrow{b})\overrightarrow{c}=\overrightarrow{a}(\overrightarrow{b}\cdot\overrightarrow{c})$}
{$(\overrightarrow{a}\cdot\overrightarrow{b})^2=\overrightarrow{a}^2\cdot\overrightarrow{b}^2$}
{$\overrightarrow{a}\cdot\overrightarrow{b}=|\overrightarrow{a}|\cdot|\overrightarrow{b}|\sin (\overrightarrow{a},\overrightarrow{b})$}
{\True $\overrightarrow{a}\cdot(\overrightarrow{b}-\overrightarrow{c})=\overrightarrow{a}\cdot\overrightarrow{b}-\overrightarrow{a}\cdot\overrightarrow{c}$}
\loigiai{
Khẳng định đúng là $\overrightarrow{a}\cdot(\overrightarrow{b}-\overrightarrow{c})=\overrightarrow{a}\cdot\overrightarrow{b}-\overrightarrow{a}\cdot\overrightarrow{c}$.
}
\end{ex}
\begin{ex}%[0H1B3-2]
Cho hình vuông $ABCD$ có cạnh $a$. Khẳng định nào sau đây là đúng?
\choice
{$(\overrightarrow{AB},\overrightarrow{BD})=45^\circ$}
{$(\overrightarrow{AC},\overrightarrow{BC})=45^\circ$ và $\overrightarrow{AC}\cdot\overrightarrow{BC}=a^2$}
{$\overrightarrow{AC}\cdot\overrightarrow{BD}=a^2\sqrt{2}$}
{$\overrightarrow{BA}\cdot\overrightarrow{BD}=-a^2$}
\loigiai{
\immini{Khẳng định đúng là $(\overrightarrow{AC},\overrightarrow{BC})=45^\circ$ và $\overrightarrow{AC}\cdot\overrightarrow{BC}=a^2$.\\
Có $(\overrightarrow{AC},\overrightarrow{BC})=(\overrightarrow{AC},\overrightarrow{AD})=45^\circ$.\\
$\overrightarrow{AC}\cdot\overrightarrow{BC}=AC\cdot BC \cdot \cos 45^\circ=a\cdot a\sqrt{2}\cdot \dfrac{\sqrt{2}}{2}=a^2$.}
{\begin{tikzpicture}[scale=1, font=\footnotesize,>=stealth]%<DTools>
%Định nghĩa điểm.
\coordinate (A) at (0,3);
\coordinate (B) at (0,0);
\coordinate (C) at (3,0);
\coordinate (D) at (3,3);
%Vẽ tứ giác ABCD.
\draw (A)--(B)--(C)--(D)--cycle;
%Hiển thị các điểm.
\foreach \x/\y in {A/90,B/180,C/0,D/90}{\fill (\x) circle(1pt) ($(\x)+(\y:0.3cm)$) node{$\x$};}
\end{tikzpicture}}
}
\end{ex}
\subsection{Bài tập tự luận}
\begin{bt}%[0H1B3-4]
Trên cạnh $BC$ của tam giác $ABC$ lấy điểm $M$ sao cho $MB=3MC$.
\begin{enumerate}
\item Tìm mối liên hệ giữa hai vec-tơ $\vec{MB}$ và $\vec{MC}$.
\item Biểu thị vec-tơ $\vec{AM}$ theo hai vec-tơ $\vec{AB}$ và $\vec{AC}$.
\end{enumerate}
\loigiai{
\immini{\begin{enumerate}
\item Có: $\vec{MB}=-3\vec{MC}$.
\item Gọi $I$ là trung điểm của $BC$. Theo quy tắc trung điểm ta có
$\vec{AI}=\dfrac{1}{2}\left(\vec{AB}+\vec{AC}\right)$.\\
Vì $M$ là trung điểm của $IC$ nên theo quy tắc trung điểm ta có\\
$\vec{AM}=\dfrac{1}{2}\left(\vec{AI}+\vec{AC}\right)=\dfrac{1}{2}\left(\dfrac{1}{2}\vec{AB}+\dfrac{3}{2}\vec{AC}\right)\\=\dfrac{1}{4}\vec{AB}+\dfrac{3}{4}\vec{AC}$.
\end{enumerate}}
{\begin{tikzpicture}[scale=.6, font=\footnotesize,>=stealth]%<DTools>
\def\canhAB{4};\def\canhBC{6};\def\gocABC{50};
%Định nghĩa điểm.
\coordinate (B) at (0,0);
\coordinate (A) at ($(B)+(\gocABC:\canhAB)$);
\coordinate (C) at ($(B)+(0:\canhBC)$);
\coordinate (M) at ($(B)!3/4!(C)$);
\coordinate (I) at ($(B)!1/2!(C)$);
%Vẽ tam giác ABC.
\draw (A)--(B)--(C)--cycle (A)--(M) (A)--(I);
%Hiển thị các điểm.
\foreach \x/\y in {A/90, B/180, C/0,M/-90,I/-90}{\fill (\x) circle(1pt) ($(\x)+(\y:0.3cm)$) node{$\x$};}
\end{tikzpicture}}
}
\end{bt}
\begin{bt}%[0H1B2-2]
Cho hình bình hành $ABC$. Chứng minh rằng với mọi điểm $M$, ta có:
$$\vec{MA}+\vec{MC}=\vec{MB}+\vec{MD}.$$
\loigiai{
\immini{Gọi $O$ là giao điểm của hai đường chéo hình bình hành.\\
Vì $O$ là trung điểm của $AC$ nên $\vec{MA}+\vec{MC}=2\vec{MO}$.\\
Vì $O$ là trung điểm của $BD$ nên $\vec{MB}+\vec{MD}=2\vec{MO}$.\\
Vậy $\vec{MA}+\vec{MC}=\vec{MB}+\vec{MD}$.}{\begin{tikzpicture}[scale=1, font=\footnotesize,>=stealth]%<DTools>
%Định nghĩa điểm.
\coordinate (A) at (1,3);
\coordinate (B) at (0,0);
\coordinate (C) at (3,0);
\coordinate (D) at (4,3);
\coordinate (O) at ($(A)!1/2!(C)$);
%Vẽ tứ giác ABCD.
\draw (A)--(B)--(C)--(D)--cycle (A)--(C) (B)--(D);
%Hiển thị các điểm.
\foreach \x/\y in {A/90,B/180,C/0,D/90,O/0}{\fill (\x) circle(1pt) ($(\x)+(\y:0.3cm)$) node{$\x$};}
\end{tikzpicture}}
}
\end{bt}

\begin{bt}%4.35
	Trong mặt phẳng tọa độ $Oxy$, cho $A(2;1)$, $B(-2;5)$ và $C(-5;2)$.
		\begin{enumerate}
			\item Tìm tọa độ của các vectơ $\overrightarrow{BA}$ và $\overrightarrow{BC}$.
			\item Chứng minh rằng $A,~B,~C$ là ba đỉnh của một tam giác vuông. Tính diện tích và chu vi của tam giác đó.
			\item Tìm tọa độ trọng tâm $G$ của tam giác $ABC$.
			\item  Tìm tọa độ của điểm $D$ sao cho tứ giác $BCAD$ là một hình bình hành.
		\end{enumerate}

	\loigiai{
			\begin{enumerate}
				\item $\overrightarrow{BA}=(4;-4)$ và $\overrightarrow{BC}=(-3;-3)$.
				\item 
				\begin{itemize}
					\item Do $\overrightarrow{BA}=(4;-4)$ và $\overrightarrow{BC}=(-3;-3)$ nên $\overrightarrow{BA}$ và $\overrightarrow{BC}$ không cùng phương hay $A,~B,~C$ là ba đỉnh của một tam giác.\\
					Mặt khác $\overrightarrow{BA}\cdot \overrightarrow{BC}=4(-3)-4(-3)=0$ nên tam giác $ABC$ vuông tại $B$.
					\item $BA=\sqrt{16+16}=4\sqrt{2}$, $BC=\sqrt{9+9}=3\sqrt{2}$, $AC=\sqrt{49+1}=5\sqrt{2}$.\\
					Diện tích của tam giác $ABC$ là $S_{ABC}=\dfrac{1}{2}\cdot BA\cdot BC=\dfrac{1}{2}\cdot 4\sqrt{2}\cdot 3\sqrt{2}=12$.\\
					Chu vi của tam giác $ABC$ là $2P=AB+AC+BC=4\sqrt{2}+5\sqrt{2}+3\sqrt{2}=12\sqrt{2}$.
				\end{itemize}
				\item Tọa độ trọng tâm $G$ của tam giác $ABC$ là $\heva{&x_G=\dfrac{x_A+x_B+x_C}{3}=-\dfrac{5}{3}\\&y_G=\dfrac{y_A+y_B+y_C}{3}=\dfrac{8}{3}} \Rightarrow G\left(-\dfrac{5}{3};\dfrac{8}{3}\right)$.
				\item Tứ giác $BCAD$ là một hình bình hành\\ $ \Leftrightarrow \overrightarrow{BC}=\overrightarrow{DA} \Leftrightarrow \heva{&x_A-x_D=-3\\&y_A-y_D=-3} \Leftrightarrow \heva{&x_D=5\\&y_D=4.}$\\
				Vậy $D=\left(5;4\right)$.
			\end{enumerate}
}
\end{bt}


\begin{bt}%4.36
	Trong mặt phẳng tọa độ $Oxy$, cho $A(1;2)$, $B(3;4)$, $C(-1;-2)$ và $D(6;5)$.

		\begin{enumerate}
			\item Tìm tọa độ của các vectơ $\overrightarrow{AB}$ và $\overrightarrow{CD}$.
			\item Hãy giải thích tại sao các vectơ $\overrightarrow{AB}$ và $\overrightarrow{CD}$ cùng phương.
			\item Giả sử $E$ là điểm có tọa độ $(a;1)$. Tìm $a$ để các vectơ $\overrightarrow{AC}$ và $\overrightarrow{BE}$ cùng phương.
			\item  Với $a$ tìm được, hãy biểu thị vectơ $\overrightarrow{AE}$ theo các véctơ $\overrightarrow{AB}$ và $\overrightarrow{AC}$.
		\end{enumerate}

	\loigiai{
			\begin{enumerate}
				\item $\overrightarrow{AB}=\left(2;2\right)$ và $\overrightarrow{CD}=\left(7;7\right)$.
				\item Do $\dfrac{2}{7}=\dfrac{2}{7}$ nên vectơ $\overrightarrow{AB}$ và $\overrightarrow{CD}$ cùng phương.
				\item Ta có $\overrightarrow{AC}=\left(-2;-4\right)$ và $\overrightarrow{BE}=\left(a-3;-3\right)$ nên vectơ $\overrightarrow{AC}$ và $\overrightarrow{BE}$ cùng phương $ \Leftrightarrow \dfrac{a-3}{-2}=\dfrac{-3}{-4} \Leftrightarrow a=\dfrac{3}{2}$.
				\item Ta có $\overrightarrow{AE}=\left(\dfrac{1}{2};-1\right)$.\\
				$\overrightarrow{AE}=m\cdot \overrightarrow{AB}+n\cdot \overrightarrow{AC} \Leftrightarrow \heva{&2m-2n=\dfrac{1}{2}\\&2m-4n=-1} \Leftrightarrow \heva{&m=1\\&n=\dfrac{3}{4}.}$\\
				Vậy 	$\overrightarrow{AE}=\overrightarrow{AB}+\dfrac{3}{4}\cdot \overrightarrow{AC}$.
			\end{enumerate}
	}
\end{bt}
\begin{bt}%4.37
	Cho vectơ $\overrightarrow{a}\ne \overrightarrow{0}$. Chứng minh rằng $\dfrac{1}{\left|\overrightarrow{a}\right|}\overrightarrow{a}$ (hay còn được viết là $\dfrac{\overrightarrow{a}}{\left|\overrightarrow{a}\right|}$) là một vectơ đơn vị cùng hướng với vectơ $\overrightarrow{a}$.
	\loigiai{Do $\overrightarrow{a}\ne \overrightarrow{0}$ nên $\left|\overrightarrow{a}\right|> 0$. Do đó $\dfrac{1}{\left|\overrightarrow{a}\right|}>0$ và $\left|\dfrac{1}{\left|\overrightarrow{a}\right|}\overrightarrow{a}\right|=\dfrac{\left|\overrightarrow{a}\right|}{\left|\overrightarrow{a}\right|}=1$.
		Suy ra $\dfrac{\overrightarrow{a}}{\left|\overrightarrow{a}\right|}$ là một vectơ đơn vị cùng hướng với vectơ $\overrightarrow{a}$.
	}
\end{bt}
\begin{bt}%4.38
	Cho ba vectơ $\overrightarrow{a},~\overrightarrow{b},~\overrightarrow{u}$ với $\left|\overrightarrow{a}\right|=\left|\overrightarrow{b}\right|=1$ và $\overrightarrow{a}\perp \overrightarrow{b}$. Xét một hệ trục $Oxy$ với các vectơ đơn vị $\overrightarrow{i}=\overrightarrow{a}$, $\overrightarrow{j}=\overrightarrow{b}$. Chứng minh rằng

		\begin{enumerate}
			\item Vectơ $\overrightarrow{u}$ có tọa độ là $\left(\overrightarrow{u}\cdot \overrightarrow{a};\overrightarrow{u}\cdot \overrightarrow{b}\right)$.
			\item $\overrightarrow{u}=\left(\overrightarrow{u}\cdot \overrightarrow{a}\right)\overrightarrow{a}+\left(\overrightarrow{u}\cdot \overrightarrow{b}\right)\overrightarrow{b}$
			
		\end{enumerate}

	\loigiai{\immini{
		
			\begin{enumerate}
				\item Ta có $\overrightarrow{u}=\overrightarrow{OH}+\overrightarrow{OK}=\overline{OH}\cdot \overrightarrow{i}+\overline{OK}\cdot \vec{j}$.\\
				Ta lại có $\vec{u}\cdot \vec{i}=\left|\vec{u}\right|\cdot \left|\vec{i}\right|\cdot \cos \left(\vec{u},\vec{i}\right)=\overline{OH}$.\\
				 $\vec{u}\cdot \vec{j}=\left|\vec{u}\right|\cdot \left|\vec{j}\right|\cdot \cos \left(\vec{u},\vec{j}\right)=\overline{OK}$.\\
				 Suy ra $\vec{u}=\left(\vec{u}\cdot \vec{i}\right)\cdot \vec{i}+\left(\vec{u}\cdot \vec{j}\right)\cdot \vec{j}=\left(\vec{u}\cdot \vec{a}\right)\cdot \vec{i}+\left(\vec{u}\cdot \vec{b}\right)\cdot \vec{j}$.\\
				 Do đó vectơ $\overrightarrow{u}$ có tọa độ là $\left(\overrightarrow{u}\cdot \overrightarrow{a};\overrightarrow{u}\cdot \overrightarrow{b}\right)$.
				\item $\vec{u}=\left(\vec{u}\cdot \vec{a}\right)\cdot \vec{i}+\left(\vec{u}\cdot \vec{b}\right)\cdot \vec{j}=\left(\vec{u}\cdot \vec{a}\right)\cdot \vec{a}+\left(\vec{u}\cdot \vec{b}\right)\cdot \vec{b}$.
			\end{enumerate}
		}
		{\begin{tikzpicture}[scale=0.7,>=stealth, font=\footnotesize, line join=round, line cap=round]
				\def\xmin{-2} \def\xmax{4}  \def\ymin{-1}  \def\ymax{4} 
				%	\draw[color=gray!50,dashed] (\xmin,\ymin) grid (\xmax,\ymax);
				\draw[->] (\xmin,0)--(\xmax,0) node [below]{$x$};
				\draw[->] (0,\ymin)--(0,\ymax) node [left]{$y$};
				\node at (0,0) [below left]{$O$};
				\draw[->] (0,0)--(1,0)node [below left]{$\overrightarrow{i}$};
				\draw[->] (0,0)--(0,1)node [below left]{$\overrightarrow{j}$};
				\draw[->] (0,0)--(3,3)node [above left]{$\overrightarrow{u}$};
				\draw[dashed] (3,0)node[below]{$H$}--(3,3)--(0,3)node[left]{$K$};
				\clip (\xmin,\ymin) rectangle (\xmax,\ymax);
		\end{tikzpicture}}
	}
\end{bt}
\begin{bt}%4.39
	Trên sông một ca nô chuyển động thẳng đều theo hướng $S15^\circ E$ với vận tốc có độ lớn bằng $20$km/h. Tính vận tốc riêng của ca nô, biết rằng nước trên sông chảy về hướng đông với vận tốc có độ lớn bằng $3$km/h.
	\loigiai{
		\immini{
			Gọi $V$ là vận tốc của ca nô chuyển động trên sông.\\
			$v_n$ là vận tốc của dòng chảy.\\
			$v_c$ là vận tốc riêng của ca nô.\\
			Ta có $\vec{v}=\vec{OE}$, $\vec{v_n}=\vec{OS}$, $\vec{v_c}=\overrightarrow{OM}$.\\
			$\vec{v}=\vec{v_c}+\vec{v_n} \Rightarrow \vec{v_c}=\vec{v}-\vec{v_n}$.\\
			$v^2_c=v^2+v^2_n-2\vec{v}\cdot \vec{v_n}=20^2+3^2-20\cdot 3\cos 15^\circ \\ \Rightarrow v_c\approx 18{,}736$ km.	}	
		{\begin{tikzpicture}[scale=.7,>=stealth, font=\footnotesize, line join=round, line cap=round]
%				\draw[color=gray!50,dashed] (-5,-5) grid (5,5);
				\draw[->] (-1,0)--(5,0) node [below]{$x$};
				\draw[->] (0,-1)--(0,3) node [left]{$y$};
				\node at (0,0) [below left]{$O$};
				\node at (0.85,0.2) [right]{$15^\circ$};
				\clip (-5,-5) rectangle (5,5);
				\draw[->] (0,0)--(15:4)circle (0.05) node[above right]{$E$};
				\draw[dashed] (0,1)node[above right]{$M$}--(15:4)--(3.8,0)node[below]{$S$};
		\end{tikzpicture}}
	}
\end{bt}

\begin{bt}%[Trong Nguyen]%[0H1K3-1]
	Chứng minh 
	\begin{enumerate}
		\item Nếu $ ABCD $ là hình bình hành thì $ \vec{AB}+\vec{AD}+\vec{CE} = 
		\vec{AE} $ với $ E $ là điểm bất kì.
		\item Nếu $ I $ là trung điểm của đoạn thẳng $ AB $ thì $ 
		\vec{MA}+\vec{MB}+2\vec{IN} = 2\vec{MN}$ với $ M $, $ N $ là hai điểm 
		bất kì.
		\item Nếu $ G $ là trọng tâm của tam giác thì $ \vec{MA}+ 
		\vec{MB}+\vec{MC}-3\vec{MN} = 3 \vec{NG}$ với $ M $, $ N $ là hai điểm 
		bất kì.
	\end{enumerate}
	\loigiai{
			\begin{enumerate}
			\item Do $ ABCD $ là hình bình hành nên $ 
			\vec{AB}+\vec{AD}=\vec{AC} $. \\
			Khi đó $ \vec{AB}+\vec{AD}+\vec{CE} = \vec{AC}+\vec{CE}=\vec{AE} $ 
			 với $ E $ là điểm bất kì.
			\item Do $ I $ là trung điểm của đoạn thẳng $ AB $ nên ta có $ 
			\vec{MA}+\vec{MB} = 2\vec{MI} $. \\
			 Khi đó $ \vec{MA}+\vec{MB}+2\vec{IN} = 2\vec{MI}+ 2\vec{IN} = 2 
			 \left( \vec{MI}+\vec{IN}\right) = 2\vec{MN}$ với $ M $, $ N $ là 
			 hai điểm bất kì.
			\item Do $ G $ là trọng tâm của tam giác nên ta có $ \vec{MA}+ 
			\vec{MB}+\vec{MC} = 3\vec{MG} $. \\
			Khi đó $ \vec{MA}+ \vec{MB}+\vec{MC}-3\vec{MN} = 3\vec{MG}- 
			3\vec{MN} = 3 \left( \vec{MG}-\vec{MN}\right) = 3\vec{NG}$ với $ M 
			$, $ N $ là hai 
			điểm 
			bất kì.
		\end{enumerate}
	
	}
	
\end{bt}


\begin{bt}%[Trong Nguyen]%[0H1Y1-3]
	Cho hình bình hành $ ABCD $ có $ AB=4 $, $ AD=6 $, $ \widehat{BAD} = 
	60^\circ $.
	\begin{enumerate}
		\item Biểu thị các vectơ $ \vec{BD} $, $ \vec{AC} $ theo $ \vec{AB} $, 
		$ \vec{AD} $.
		\item Tính các tích vô hướng $ \vec{AB} \cdot \vec{AD} $, $ \vec{AB} 
		\cdot \vec{AC}$, $ \vec{BD} \cdot \vec{AC}$.
		\item Tính độ dài các đường chéo $ BD $, $ AC $.
	\end{enumerate}
	\loigiai{
		\immini
		{
			a) Ta có $ \vec{BD} = \vec{AD}-\vec{AB} $. \\
				$ \vec{AC} = \vec{AB}+\vec{AD} $.
			
		}
		{
		\begin{tikzpicture}[>=stealth,scale=.6, line join = round, line cap 
		= round]
		\coordinate (A) at (0,0);
		\coordinate (B) at (60:4);
		\coordinate (D) at (6,0);
		\coordinate (C) at ($(B)-(A)+(D)$);
		\draw[blue,line width=1] (A) node[left] {$ A $}--node[left]{$ 4 $} 
		(B)node[left] {$ 
		B $}--
		(C)node[right] {$ C $}--(D)node[right] {$ D $}--node[below]{$ 6 $} 
		cycle;
		\pic
		[draw, "$60^\circ$", angle eccentricity=1.8] {angle = D--A--B};
		\end{tikzpicture}
		}
\noindent	b) 	$ \vec{AB} \cdot \vec{AD} = \left|\vec{AB}\right| \cdot 
	\left|\vec{AD}\right| \cos \left(\vec{AB}, \vec{AD}\right) = 4 
	\cdot 6 \cdot \cos\widehat{BAD} = 24 \cdot \cos 60^\circ = 24 
	\cdot \dfrac{1}{2} = 12  $. \\
	$ \vec{AB} \cdot \vec{AC} = \vec{AB} \cdot \left( \vec{AB} + 
	\vec{AD}\right) = \vec{AB} \cdot \vec{AB}+ \vec{AB} \cdot 
	\vec{AD} = 4^2+12=28   $. \\
	$ \vec{BD} \cdot \vec{AC} = \left( \vec{AD} - 
	\vec{AB}\right) \cdot \left( \vec{AB} + 
	\vec{AD}\right) = \left( \vec{AD} -	\vec{AB}\right) \cdot \vec{AB} + \left( 
	\vec{AD} -	\vec{AB}\right) \cdot \vec{AD} =  \vec{AD} \cdot \vec{AB}  
	-\vec{AB} \cdot \vec{AB} + \vec{AD} \cdot \vec{AD}- \vec{AB} 
	\cdot \vec{AD} = - \vec{AB}^2 + \vec{AD}^2 = -4^2+6^2 = 20 $.\\
	c) Ta có $ \vec{BD}^2= \left( \vec{AD} - \vec{AB}\right)^2 = \vec{AD}^2 + 
	\vec{AB}^2 -2 \vec{AB} \cdot \vec{AD} = 6^2+4^2-2 \cdot 12 = 28  $. \\
	Suy ra $ BD^2=28 \Rightarrow BD = 2\sqrt{7} $. \\
	Ta có $ \vec{AC}^2= \left( \vec{AB} + \vec{AD}\right)^2 = \vec{AB}^2 + 
	\vec{AD}^2 	+2 \vec{AB} \cdot \vec{AD} = 4^2+6^2+2 \cdot 12 = 66  $. \\
	Suy ra $ AC^2=66 \Rightarrow AC=\sqrt{66} $.	
	}	
\end{bt}

\begin{bt}%[Trong Nguyen]%[0H1K3-1]
	Hai lực $ \vec{F_1} $, $ \vec{F_2} $ cho trước cùng tác dụng lên một vật 
	tại điểm $ O $ và tạo với nhau một góc $ \left( \vec{F_1}, \vec{F_2}\right) 
	= \alpha $ làm cho vật di chuyển từ $ O $ đến $ C $. Lập công thức tính 
	cường độ của hợp lực $ \vec{F} $ làm cho vật di chuyển theo hướng từ $ O $ 
	đến $ C $ (giả sử chỉ có đúng hai lực $ \vec{F_1} $, $ \vec{F_2} $ làm cho 
	vật di chuyển).
	\loigiai{
		\immini
		{
		Ta có $ \vec{F}^2= \left( \vec{F_1} + \vec{F_2}\right)^2 = \vec{F_1}^2 
		+ \vec{F_2}^2 +2 \vec{F_1} \cdot \vec{F_2} = \vec{F_1}^2 + 
		\vec{F_2}^2+2 \left| \vec{F_1} \right| \cdot \left| \vec{F_2} \right| 
		\cdot \cos \left( \vec{F_1}, \vec{F_2}\right) = \left| \vec{F_1} 
		\right|^2 + \left| \vec{F_1} \right|^2+2 \left| \vec{F_2} \right| \cdot 
		\left| \vec{F_2} \right| 
		\cdot \cos \alpha    $. \\
		Suy ra $ \left| \vec{F} \right| = \sqrt{\left| \vec{F_1} \right|^2 + 
			\left| \vec{F_2} \right|^2+2 \left| \vec{F_1} \right| \cdot \left| 
			\vec{F_2} 
			\right| \cdot \cos \alpha } $.
		}
		{
			\begin{tikzpicture}[>=stealth,scale=.6, line join = round, line cap 
			= round]
			\coordinate[label=left:$ O $] (O) at (0,0);
			\coordinate[label=left:$ A $] (A) at (60:4);
			\coordinate[label=right:$ B $] (B) at (6,0);
			\coordinate[label=right:$ C $] (C) at ($(B)-(O)+(A)$);
			\draw[->,blue,line width=1] (O)--node[left] {$ \vec{F_1} $}(A);
			\draw[->,blue,line width=1](O)--node[below]{$ \vec{F_2} $} (B);
			\draw[->,blue,line width=1] (O)--node[above left]{$ \vec{F} $} (C);
			\draw[dashed] (A)--(C)--(B);
			\pic
			[draw, "$\alpha$", angle eccentricity=1.5] {angle = B--O--A};
			\end{tikzpicture}
		}
	}	
\end{bt}

\Closesolutionfile{ans}